\documentclass[spanish,letterpaper,oneside]{article}
\def\documenttitle {Delitos de peligro efectivo y presunto}
\def\documentsubtitle {Riesgo, cuidado y responsabilidad penal}
\def\documentsubject {Actividad 4 - Derecho penal especial mexicano}
\def\documentauthor {Martín Jonathan de la Cruz}
\def\coursename {Derecho penal especial mexicano}
\def\coursecode {020}
\def\universityname {Universidad Abierta y a Distancia de México}
\def\universityfaculty {Licenciatura en Derecho}
\def\universitydepartment {Derecho penal especial mexicano}
\def\universitydepartmentimage {departamentos/UnADM}
\def\universitydepartmentimagecfg {height=1.57cm}
\def\universitylocation {Roma Norte, Ciudad de México}
\def\authortable {
  \begin{tabular}{ll}
    \textbf{Alumno:} & Martín Jonathan de la Cruz \\
    Matrícula: & ES2611202040 \\
    Figura docente: & Verónica Ortiz López \\
    Grupo: & 020 \\
    Semestre/Bloque: & 2 / 2 \\
    Semana: & 4 \\
    \multicolumn{2}{l}{Fecha: 23 de octubre de 2026} \\
    \multicolumn{2}{l}{\universitylocation}
  \end{tabular}
}
\input{template}
\def\indexstyle{}
\usepackage{helvet}
\renewcommand{\familydefault}{\sfdefault}
\usepackage[most]{tcolorbox}
\usepackage{tabularx}
\usepackage{tikz}
\usetikzlibrary{babel}
\definecolor{casoverde}{HTML}{174A3A}
\definecolor{casopapel}{HTML}{F6F7F2}
\setcitestyle{authoryear,open={(},close={)},citesep={;}}
\def\pagemargintop{2.3}
\def\pagemarginbottom{2.3}
\def\predocpageromannumber{false}
\newcommand{\insertcoverwatermark}{%
  \AddToShipoutPictureBG*{%
    \begin{tikzpicture}[remember picture,overlay]
      \node[opacity=0.13,inner sep=0pt,yshift=-0.35cm]
        at (current page.center) {\includegraphics[width=0.72\paperwidth]{img/departamentos/UnADM.pdf}};
    \end{tikzpicture}%
  }%
}
\newtcolorbox{casoestudio}[1]{
  enhanced,colback=casopapel,colframe=casoverde,
  boxrule=0pt,leftrule=2.4pt,arc=1mm,
  left=2.4mm,right=2.4mm,top=1.7mm,bottom=1.7mm,
  title={#1},colbacktitle=casoverde,coltitle=white,
  fonttitle=\bfseries,fontupper=\fontsize{10.5}{11.8}\selectfont,
  before skip=5pt,after skip=5pt
}
\newcommand{\fichacaso}[1]{%
  \begin{tcolorbox}[colback=white,colframe=casoverde!25,
    boxrule=0.4pt,arc=0.6mm,left=1.6mm,right=1.6mm,top=1mm,bottom=1mm,
    before skip=0pt,after skip=3pt]
    #1
  \end{tcolorbox}%
}
\setlength{\parindent}{0pt}
\setlength{\parskip}{3pt}
\setlength{\tabcolsep}{4pt}
\renewcommand{\arraystretch}{1.08}
\begin{document}
\insertcoverwatermark
\templatePortrait
\templatePagecfg
\templateFinalcfg
\hypersetup{urlcolor=casoverde,citecolor=black,linkcolor=casoverde}
\pagenumbering{arabic}\setcounter{page}{2}
\fontsize{10.5}{11.8}\selectfont
\sectionfont{\color{casoverde}\fontsize{14}{16}\selectfont\bfseries}
\subsectionfont{\color{casoverde}\fontsize{11.3}{13}\selectfont\bfseries}
\titlespacing*{\section}{0pt}{7pt}{4pt}
\titlespacing*{\subsection}{0pt}{5pt}{3pt}
\section{Introducción}
Disparar sin herir, dejar a alguien sin cuidado o exponerlo a contagio plantea un límite: prevenir daños sin castigar por intuición. Por ello, la legalidad exige una ``ley exactamente aplicable al delito de que se trata'' (CPEUM, art.~14) \citep{cpeumS4Vigente}. Como los casos no indican lugar ni fecha, los códigos federal (CPF), de Ciudad de México (CPDF) y de Nuevo León (CPNL) permiten contrastar hipótesis; su aplicación depende de competencia y vigencia, no de elegir la sanción más severa.

\section{El riesgo penal y los deberes de protección}
\subsection*{Elementos del peligro y regulación de las conductas}
La \textbf{tipicidad} exige que los hechos satisfagan todos los elementos del tipo; el \textbf{bien jurídico tutelado} es el interés protegido. El \textbf{peligro efectivo} exige exposición concreta; el \textbf{peligro presunto} anticipa la tutela frente a conductas legalmente consideradas peligrosas. Ninguno exige lesión. Sin embargo, la clasificación no elimina los requisitos de peligro expresamente previstos \citep[p.~279]{amuchateguiPenal2012}. La \textbf{omisión} incumple una acción debida y posible; la \textbf{posición de garante} es un deber especial de protección, distinto de la obligación ocasional del observador (CPDF, art.~16) \citep{cpdfS4Vigente}.

% Técnica: valoración de decisiones. Organización: conducta, norma y aplicación a hechos.
% Síntesis y juicio: distinguir elementos comprobados, alternativas y diligencias pertinentes.
\begingroup\fontsize{9.5}{10.5}\selectfont
\begin{tabularx}{\linewidth}{@{}p{3.1cm}X p{4.15cm}@{}}
\toprule
\textbf{Figura} & \textbf{Conducta y condición decisiva} & \textbf{Artículo y legislación} \\
\midrule
Disparo sobre persona & Accionar un arma sin causa justificada; el tipo local no exige lesión ni disparar contra alguien. & CPDF 211 Ter. CPF 306 derogado. \\
Ataque peligroso & Atacar con capacidad de causar muerte por el arma, fuerza, destreza u otra circunstancia. & CPNL 332 y 334; no figura autónoma en CPF ni CPDF revisados. \\
Abandono de niños y enfermos & Abandonar a niño incapaz de cuidarse o persona enferma, teniendo obligación de cuidado. & CPF 335; CPDF 156 (incapaz de valerse por sí). \\
Abandono de atropellados & Tras atropellar culposa o fortuitamente, omitir auxilio o asistencia pudiendo prestarlos. & CPF 341; CPDF 157 comprende lesión previa. \\
Abandono de cónyuge o hijos & Abandonar sin justificación y sin recursos de subsistencia; en CDMX, incumplir alimentos. & CPF 336; equivalente local: CPDF 193; insolvencia deliberada, 194. \\
Omisión de socorro / cuidado & Socorro: no avisar o auxiliar a persona abandonada en peligro, sin riesgo propio. Cuidado: abandono por obligado. & CPF 340 / 335; CPDF 156 para cuidado, sin identidad con socorro general. \\
Peligro de contagio & Sabiéndose enfermo en el supuesto legal y período infectante, exponer mediante un medio transmisible. & CPF 199 Bis, vigente; CPDF 159 derogado desde 2024. \\
\bottomrule
\end{tabularx}
\par\vspace{2pt}\textit{Fuentes:} textos oficiales CPF, CPDF y CPNL \citep{cpfS4Vigente,cpdfS4Vigente,cpnlS4Vigente}. La derogación del CPDF 159 corrige la referencia aún incluida en el material docente \citep[p.~8]{unadmPenalS4Material}.
\endgroup

\begin{casoestudio}{Caso 1. Disparo de Roberto y agresión de Miguel}
\fichacaso{\textbf{Hechos y problema.} Roberto dispara un arma de fuego hacia Carlos sin lesionarlo; Miguel intenta agredir con un objeto peligroso, pero lo detienen antes de alcanzarlo. ¿Existe delito aunque no haya daño?}
\textbf{Análisis.} Para Roberto, el CPDF 211 Ter permite investigar el disparo injustificado; comprobar arma, autoría, trayectoria, distancia y una posible causa de justificación. La tentativa de homicidio exige además intención de matar, actos ejecutivos idóneos y falta de consumación ajena a su voluntad (CPDF, arts.~20 y 123); no surge automáticamente de disparar. En cambio, para Miguel, el CPNL 332 exige un ataque con capacidad mortal; el objeto por sí solo no lo acredita. Reconstruir movimientos, zona corporal, fuerza y momento de la intervención: si ya ejecutó el ataque, puede consumarse sin impacto; si no, distinguir preparación de tentativa (CPNL, art.~31). El CPF 306 derogado no fundamenta una imputación actual. \textbf{Conclusión:} ambas hipótesis requieren prueba y ley territorialmente aplicable \citep{cpdfS4Vigente,cpnlS4Vigente,cpfS4Vigente}.
\end{casoestudio}

\clearpage
\subsection*{Abandono, auxilio y exposición al contagio}
\begin{casoestudio}{Caso 2. Ana, el conductor y Pedro}
\fichacaso{\textbf{Hechos y problema.} Ana deja solo varias horas al menor a su cargo; un conductor atropella y se retira sin pedir ayuda; Pedro omite alimentos indispensables para sus hijos. Son tres hechos independientes.}
\textbf{Análisis.} La protección de vida, salud y subsistencia depende de deberes diferentes. Para decidir fundadamente, comprobar:
\begin{itemize}[leftmargin=*,topsep=1pt,itemsep=1pt,parsep=0pt]
\item \textbf{Ana:} edad, autonomía del menor, condiciones de desamparo y obligación de cuidado; entrevistar y documentar la custodia. Varias horas sin compañía no bastan por sí solas para abandono (CPF 335; CPDF 156).
\item \textbf{Conductor:} atropellamiento culposo o fortuito, necesidad de asistencia y posibilidad efectiva de auxiliar o pedirla; revisar video, llamadas y peritaje vial. CPF 341 distingue la omisión del atropellamiento; CPDF 157 exige lesión previa.
\item \textbf{Pedro:} filiación, obligación alimentaria, recursos, pagos e impedimentos reales. CPF 336 exige abandono injustificado sin subsistencia; CPDF 193 sanciona incumplimiento alimentario aun si terceros cubren las necesidades. No consta insolvencia deliberada para aplicar CPDF 194.
\end{itemize}
\textbf{Conclusión:} investigar cada omisión y atender a las personas dependientes, sin equiparar incapacidad económica con evasión voluntaria ni sumar automáticamente delitos \citep{cpfS4Vigente,cpdfS4Vigente}.
\end{casoestudio}

\begin{casoestudio}{Caso 3. Socorro, cuidado y contagio}
\fichacaso{\textbf{Hechos y problema.} Un observador no pide ayuda; una cuidadora abandona a un enfermo dependiente; otra persona conoce una enfermedad transmisible y realiza una conducta de posible riesgo.}
\textbf{Análisis.} La \textbf{omisión de socorro} (CPF 340) exige encontrar abandonada a una persona en las condiciones del tipo, conocer su necesidad y poder avisar de inmediato a la autoridad o auxiliar sin peligro personal; observar cualquier riesgo no basta. La \textbf{omisión de cuidado} presupone obligación previa: respecto de la cuidadora, acreditar encargo aceptado, dependencia, abandono y posibilidad de actuar (CPF 335; CPDF 156). El CPDF 157 no se aplica al mero espectador: exige haber lesionado previamente culposa o fortuitamente.

El \textbf{peligro de contagio} del CPF 199 Bis exige mal venéreo u otra enfermedad grave en período infectante, conocimiento y exposición por relaciones sexuales u otro medio transmisible. Solicitar peritaje médico sobre diagnóstico, infectividad y vía concreta; no basta el diagnóstico ni se exige infección consumada. En CDMX, el art.~159 está derogado: no puede recuperarse por analogía. \textbf{Conclusión:} por tanto, diferenciar ayuda ocasional, cuidado debido y riesgo de transmisión; comprobar cada presupuesto y la competencia antes de atribuir delito \citep{cpfS4Vigente,cpdfS4Vigente}.
\end{casoestudio}

\subsection*{Decisión fundamentada}
\begin{casoestudio}{Decisión sobre el caso 2: protección del menor y prueba del abandono}
Considero prioritario asegurar el cuidado del menor por el ``interés superior de la niñez'' (CPEUM, art.~4) e investigar a Ana por posible abandono (CPDF 156), si fuera competente CDMX. Confirmaría custodia, edad, autonomía y desamparo con documentos y testimonios. La jurisprudencia 1a./J.~26/2014 exige ``pruebas de cargo suficientes'' para condenar: protegería al menor sin anticipar culpabilidad; la calificación seguiría sujeta a esas pruebas \citep{cpeumS4Vigente,cpdfS4Vigente,scjnPresuncionS4}.
\end{casoestudio}

\clearpage
\section{Conclusiones}
La ausencia de lesión no vuelve lícita una conducta peligrosa, pero tampoco reemplaza los elementos del delito. Sostengo que proteger a quien depende de otros exige actuar ante el riesgo y probar la responsabilidad por separado. La Constitución manda que ``se presuma su inocencia'' (art.~20, B, I): ni una sospecha ni un artículo derogado justifican castigar \citep{cpeumS4Vigente}.\footnote{Se utilizó inteligencia artificial (Codex de OpenAI) para organizar fuentes, contrastar normas y jurisprudencia y elaborar la redacción y maquetación. El estudiante debe revisar y asumir el análisis antes de su entrega.}

\begingroup\fontsize{10}{11.5}\selectfont
\setlength{\bibsep}{4pt}
\begin{thebibliography}{8}
\bibitem[Amuchategui Requena(2012)]{amuchateguiPenal2012}
Amuchategui Requena, I.~G. (2012). \emph{Derecho penal} (4.ª ed.). Oxford University Press México.

\bibitem[Cámara de Diputados(2026a)]{cpfS4Vigente}
Cámara de Diputados del H. Congreso de la Unión. (2026a). \emph{Código Penal Federal} (reforma DOF 13 de marzo de 2026). \url{https://www.diputados.gob.mx/LeyesBiblio/pdf/CPF.pdf}

\bibitem[Cámara de Diputados(2026b)]{cpeumS4Vigente}
Cámara de Diputados del H. Congreso de la Unión. (2026b). \emph{Constitución Política de los Estados Unidos Mexicanos} (reforma DOF 7 de octubre de 2026). \url{https://www.diputados.gob.mx/LeyesBiblio/pdf/CPEUM.pdf}

\bibitem[Congreso de la Ciudad de México(2026)]{cpdfS4Vigente}
Congreso de la Ciudad de México. (2026). \emph{Código Penal para el Distrito Federal} (reforma GOCDMX 26 de agosto de 2026). \url{https://www.congresocdmx.gob.mx/archivo-deae456e36ecb11eac420f0f1839cc70e9ea4685.pdf}

\bibitem[Congreso de Nuevo León(2026)]{cpnlS4Vigente}
Congreso del Estado de Nuevo León. (2026). \emph{Código Penal para el Estado de Nuevo León} (reforma PO 15 de junio de 2026). \url{https://www.hcnl.gob.mx/trabajo_legislativo/leyes/codigos/codigo_penal_para_el_estado_de_nuevo_leon/}

\bibitem[Suprema Corte de Justicia de la Nación(2014)]{scjnPresuncionS4}
Suprema Corte de Justicia de la Nación. (2014). \emph{Presunción de inocencia como estándar de prueba} [Jurisprudencia 1a./J. 26/2014 (10a.), Primera Sala; registro digital 2006091]. \emph{Gaceta del Semanario Judicial de la Federación}, libro 5, tomo I, 476. \url{https://sjf2.scjn.gob.mx/detalle/tesis/2006091}

\bibitem[UnADM(2026)]{unadmPenalS4Material}
Universidad Abierta y a Distancia de México. (2026). \emph{Material de apoyo S4: Delitos de peligro efectivo y presunto} [Material de curso]. Aula virtual. \url{https://aulavirtual.unadmexico.mx/mod/resource/view.php?id=217986}
\end{thebibliography}
\endgroup
\end{document}
